\documentclass{article}

\usepackage[margin=1in]{geometry}
\usepackage{graphicx}
\usepackage{float}
\usepackage{enumitem}
\usepackage[version=4]{mhchem}
\usepackage{siunitx}
\usepackage{hyperref}

\begin{document}

\begin{center}
  {\Large LATEX Workshop\par}
  \vspace{0.5em}
  {\large Exercises\par}
  \vspace{1.5em}
  \today
\end{center}

\bigskip

\section*{Install instructions}
Overleaf is an online LaTeX editor, which can be found on:

\begin{center}
  \url{www.overleaf.com}
\end{center}

First, create a new project on Overleaf in a new tab. Alternatively, copy this project (File $\rightarrow$ Make a copy) and create a new \texttt{.tex} file.

\section{Preamble}
In \LaTeX, the packages you use and document-wide settings are placed in the preamble.

\begin{enumerate}[label=\alph*)]
  \item Add the document class article.
  \item Add the package graphicx.
  \item Add metadata for a title, author and date (the command \verb|\today| gets the current date).
\end{enumerate}

\section{Content}
\begin{enumerate}[label=\alph*)]
  \item Add the document environment.
  \item Add a title to the document.
  \item Create a numbered section named ``Formatting''.
  \item Go to the link \url{http://randomtextgenerator.com/} and copy the random text in your section.
  \item Format atleast one sentence to bold, one sentence to italic and change the size of two sentences.
\end{enumerate}

\newpage

\section{Lists and references}
\begin{enumerate}[label=\alph*)]
  \item Create a section named ``Lists''.
  \item Create a numbered list and write a sentence in the first item.
  \item After the list add a reference to the first item made in b).
\end{enumerate}

\section{Floats}
\begin{enumerate}[label=\alph*)]
  \item Create a new section named ``Floats''
  \item Create an image float.
  \item Add a caption to the image.
  \item Add a label to the image.
  \item Restrict the float to the Floats section.
\end{enumerate}

\section{Tables}
Create a table with containing a finished tic-tac-toe with a label and caption.

\section{Chemistry in TeX}
Chemistry writing needs consistent notation for formulas, reactions, units, and references.

\begin{enumerate}[label=\alph*)]
  \item Create a new section named ``Chemistry in TeX''.
  \item Add the package ``mhchem''.
  \item Add the package ``siunitx''.
  \item Create chemical notation for a formula, an ion, and a reaction using \verb|\ce|.
  \item Create quantities with units using \verb|\SI| and a number using \verb|\num|.
  \item Put a reaction scheme, spectrum, or table in a float with a caption and label so it can be referenced later.
\end{enumerate}

\newpage

\section{BibTeX}
\begin{enumerate}
  \item Go to \url{http://scholar.google.com} and search for any paper, e.g. ``BibTeX''.
  \item Obtain at least two BibTeX references and save them into one file under the .bib extension.
  \item Cite your references using the numeric style, which can be done using unsrt in bibliography style. For an example, see \url{www.bibtex.org/using}
\end{enumerate}

\end{document}

-also text down here is not shown in the pdf :)
